\newcommand{\tipoRequerimento}{cobrança indevida por perda nos reatores de lâmpadas}

% Definição de condicionais (O script deve definir como true ou false)
\newif\ifperdaVapor
\perdaVaporfalse

\newif\ifperdaFluorescente
\perdaFluorescentetrue

\section{DOS FATOS E DOS DIREITOS}

\begin{enumerate}[resume]
\item 
Trata-se de reclamação administrativa interposta com vistas à restituição de valores indevidamente faturados devido a cobrança indevida realizada pela Concessionária, consistente na inclusão de valores referentes a “perdas em reatores” de lâmpadas fluorescentes no faturamento do parque de iluminação pública do Município, apurado exclusivamente por estimativa.

\item 
Nos termos do art. 473 da Resolução Normativa ANEEL nº 1.000/2021, o faturamento da energia elétrica consumida por equipamentos auxiliares deve observar, no máximo, os limites de perdas estabelecidos nas normas técnicas aplicáveis, aferidas em condição “a frio”, bem como as especificações metrológicas expedidas pelo INMETRO, conforme Portaria nº 35/2022.

\ifperdaVapor
\item % VAPOR
Para os reatores de lâmpadas , os limites máximos admissíveis de perdas devem observar as normas técnicas aplicáveis da ABNT, as quais possuem caráter vinculante e de observância obrigatória pela Concessionária.
\fi

\ifperdaFluorescente
\item % FLUORESCENTE
No que se refere às lâmpadas fluorescentes, inexiste norma técnica da ABNT ou previsão regulatória que autorize a imputação de perdas adicionais de reator para fins de faturamento. Ademais, no caso específico das lâmpadas fluorescentes compactas, o reator integra o próprio conjunto do equipamento, razão pela qual a potência nominal declarada já contempla eventuais perdas inerentes ao seu funcionamento. Assim, a inclusão de valores a título de “perdas de reator” no faturamento revela-se desprovida de base técnica, normativa e jurídica, configurando cobrança indevida, em afronta direta ao art. 473 da REN nº 1.000/2021 e aos princípios da legalidade e da modicidade tarifária. 
\fi

\item 
Tal entendimento encontra-se pacificado e reiteradamente consolidado no âmbito da ANEEL, conforme consignado no Despacho nº 1.456/2023-STD/ANEEL, bem como em sucessivas decisões proferidas por sua Diretoria Colegiada nos processos:
\begin{itemize}
\item nº 48500.002048/2023-42, 
\item nº 48500.002050/2023-11, 
\item nº 48500.006919/2022-16, 
\item nº 48500.004731/2023-14, 
\item nº 48500.005659/2023-42, 
\item nº 48500.002136/2023-44,
\item nº 48500.006249/2023-19,
\item nº 48500.006121/2023-55.
\end{itemize}

\item 
Em todos os precedentes listados anteriormente, a Agência Reguladora determinou a estrita observância dos limites máximos de perdas “a frio” definidos pelas normas técnicas da ABNT, vedando expressamente a realização de cobranças baseadas em estimativas genéricas ou desprovidas de respaldo técnico-normativo, sob pena de caracterização de cobrança indevida.

\item
À luz do entendimento reiterado e uniforme firmado pela ANEEL nos julgamentos supracitados, destaca-se, a título exemplificativo, o posicionamento adotado no julgamento do Processo nº 48500.006121/2023-55, no qual restou expressamente reconhecida a ilegalidade da inclusão de perdas não comprovadas e não autorizadas normativamente no faturamento da classe Iluminação Pública, reafirmando-se o dever da Concessionária de limitar a cobrança aos parâmetros técnicos reguladoramente admitidos, bem como de proceder à restituição dos valores faturados a maior:

\Citacao{DESPACHO Nº 320, DE 11 DE FEVEREIRO DE 2025 

[...] 

DECIDE:  

(i) conhecer do recurso interposto pela Enel Distribuição Ceará, inscrita no CNPJ sob o nº 07.047.251/0001-70 em face da decisão emitida pela Diretoria da Agência Reguladora do Estado do Ceará – ARCE no processo nº PROC/OUV/PROC 11411/2022, para, no mérito, negar-lhe provimento;  

(ii) conhecer do recurso interposto pelo Município de Jaguaribe – CE, inscrita no CNPJ sob o nº 07.443.708/0001-66 em face da decisão emitida pela Diretoria da Agência Reguladora do Estado do Ceará – ARCE no processo nº PROC/OUV/PROC 11411/2022, para, no mérito, dar-lhe parcial provimento no sentido de:  

(ii.a) reformar parcialmente a decisão exarada pelo Conselho Diretor da ARCE no Processo PROC/OUV/PROC 11411/2022 e determinar que a Enel Distribuição Ceará corrija as perdas aplicadas no faturamento do sistema de IP do Município de Jaguaribe – CE de forma a contemplar, no cálculo do consumo dos equipamentos auxiliares, as alterações das normas ABNT referentes a cada tipo de lâmpada, aplicando-se as perdas a frio estabelecidas nas normas ABNT;  

(ii.b) determinar que a Enel Distribuição Ceará revise os faturamentos do sistema de IP do Município de Jaguaribe – CE de forma a contemplar, no cálculo do consumo dos equipamentos auxiliares referentes a todas as lâmpadas de vapor de sódio, as alterações das normas ABNT referentes a cada tipo de lâmpada, aplicandose as perdas a frio estabelecidas nas normas ABNT, seguindo os procedimentos dispostos no art. 113 da Resolução Normativa nº 414, de 2010, observando-se o Despacho ANEEL nº 18, de 2019, pelo período de 30/07/2011 até a data da efetiva correção dos valores de perdas do faturamento, com a devolução realizada em dobro, podendo abater do total a devolver os valores já efetivamente devolvidos ao Município;  

(ii.c) determinar que a Enel Distribuição Ceará revise os faturamentos do sistema de IP do Município de Jaguaribe – CE de forma a excluir, no cálculo do consumo dos equipamentos auxiliares referentes a todas as lâmpadas fluorescentes, as perdas consideradas, visto que a perda indicada pela ABNT é zero, seguindo os procedimentos dispostos no art. 113 da Resolução Normativa nº 414, de 2010, observando-se o Despacho ANEEL nº 18, de 2019, pelo período de 30/07/2011 até a data da efetiva correção dos valores de perdas do faturamento, com a devolução realizada em dobro, podendo abater do total a devolver os valores já efetivamente devolvidos ao Município;  

(ii.d) determinar à Enel Distribuição Ceará enviar aos representantes do Município o detalhamento dos cálculos dos valores devolvidos, conforme art. 133 da Resolução Normativa nº 414, de 2010, discriminando os valores faturados incorretamente para cada tipo de lâmpada, atualização e juros incidentes;  

(iii) determinar que esta decisão seja cumprida no prazo de 15 (quinze) dias após o seu trânsito em julgado; e  

(iv) determinar que a distribuidora envie à ANEEL, num prazo máximo de 15 (quinze) dias após o prazo previsto no item

(viii) desta decisão, comprovação do seu cumprimento. (grifou-se)}

\item
Diante do arcabouço normativo e do entendimento consolidado no âmbito da ANEEL, o tratamento regulatório aplicável à matéria impõe as seguintes balizas: 

\begin{enumerate}
    \ifperdaVapor
    \item % VAPOR
    Para reatores de lâmpadas , admite-se, no máximo, a consideração das perdas estritamente previstas nas respectivas normas técnicas da ABNT, aferidas nos termos regulamentares.
    \fi
    \ifperdaFluorescente
    \item % FLUORESCENTE
    Para lâmpadas fluorescentes, diante da inexistência de norma técnica específica que autorize a imputação de perdas adicionais  e, no caso das fluorescentes compactas, por se tratar de conjunto indissociável entre lâmpada e reator, as perdas adicionais devem ser fixadas em zero, sendo vedada qualquer majoração por estimativa.
    \fi
\end{enumerate}

\ifperdaVapor
\item % VAPOR
No caso concreto, verifica-se que os percentuais de perdas adotados pela Concessionária extrapolam os limites máximos admitidos pelas normas técnicas aplicáveis, configurando cobrança indevida e prática incompatível com o regime regulatório vigente. Tal conduta afronta diretamente o art. 473 da REN nº 1.000/2021, o princípio da legalidade administrativa (art. 37, caput, da Constituição Federal) e as disposições do Código de Defesa do Consumidor, notadamente os arts. 6º, inciso III, 39, inciso V, e 42, parágrafo único. Impõe-se, portanto, a devolução dos valores faturados a maior, nos termos do art. 323, inciso II, da REN nº 1.000/2021, em consonância com o entendimento firmado no Despacho nº 2.006/2024, acrescida dos devidos encargos legais. 

\item % VAPOR
Constata-se, ainda, que os percentuais de perdas considerados no faturamento realizado pela Concessionária ultrapassam os limites máximos previstos nas normas técnicas da ABNT aplicáveis, em manifesta desconformidade com o art. 473 da REN nº 1.000/2021 e com o entendimento reiterado e vinculante da ANEEL, conforme demonstrado de forma objetiva na tabela comparativa apresentada a seguir. 

 % VAPOR

\item % VAPOR
Registre-se que a Resolução Normativa ANEEL nº 1.000/2021 consolidou, em linhas 
gerais, o entendimento já firmado pela revogada Resolução Normativa nº 414/2010, 
mantendo a diretriz regulatória quanto aos critérios de faturamento da classe 
Iluminação Pública. Nesse sentido, dispõe o art. 473 da REN nº 1.000/2021: 

\Citacao{Art. 473. Para fins de faturamento, a energia elétrica consumida pelos equipamentos auxiliares de iluminação pública deve ser estimada pelos critérios das normas vigentes da Associação Brasileira de Normas Técnicas – ABNT. Parágrafo único. Mediante acordo prévio entre a distribuidora e o poder público municipal, a estimativa disposta no caput pode ser realizada por meio de dados do fabricante dos equipamentos ou em ensaios realizados em laboratórios acreditados por órgão oficial.} % VAPOR
\fi

\ifperdaVapor
\item % VAPOR
Destaca-se que o dispositivo não autoriza a adoção unilateral de percentuais genéricos ou estimativas dissociadas das normas técnicas aplicáveis, tampouco a imputação de perdas não comprovadas ou não previstas em norma específica, reforçando a ilegalidade da cobrança realizada pela Concessionária no caso concreto.
\fi

\ifperdaFluorescente
\item % FLUORESCENTE
No que se refere às lâmpadas fluorescentes, é expressamente vedada a cobrança de valores a título de “perdas de reator”, haja vista a inexistência de respaldo em norma técnica da ABNT que autorize tal imputação e, especialmente, porque, no caso das lâmpadas fluorescentes compactas, o reator integra o próprio conjunto do equipamento, estando eventuais perdas já incorporadas à potência nominal declarada pelo fabricante. 

\item % FLUORESCENTE
Não obstante essa vedação técnica e regulatória, verifica-se que a Concessionária procedeu à cobrança indevida de perdas em reatores nos faturamentos mensais da unidade consumidora da classe Iluminação Pública, realizados por estimativa e sem medição efetiva, em manifesta afronta ao art. 473 da Resolução Normativa ANEEL nº 1.000/2021 e ao entendimento reiterado e consolidado da ANEEL, conforme evidenciado na tabela demonstrativa apresentada a seguir.

\begin{figure}[H]
    \centering
    \includegraphics[width=\textwidth]{C:/Users/Usuário 1/OneDrive/PASTA ENERGIA 1/PARCEIROS/Municípios Thamires e Ruda/Paraíba/Salgado de São Félix/Reclamações/REC-006_2026-SALGADO_DE_SAO_FELIX-PERDA_NOS-REATORES/PERDA NOS REATORES - SALGADO DE SÃO FÉLIX.png}
\end{figure} % FLUORESCENTE
\fi

\item
Os dados e quantitativos acima apresentados foram extraídos do Quadro de Iluminação Pública – QIP, referente a Julho de 2021, em anexo, documento elaborado e utilizado pela própria Concessionária como base de cálculo para o faturamento mensal da unidade consumidora da classe Iluminação Pública. Ressalta-se que referido faturamento é realizado por estimativa, a partir das informações constantes no QIP, circunstância que reforça a necessidade de estrita observância dos limites técnicos e regulatórios aplicáveis, bem como evidencia a responsabilidade exclusiva da Concessionária pela correção dos parâmetros adotados.

\item
Diante do exposto, apura-se a ocorrência de cobrança indevida correspondente a 423,12 kWh/mês, relativa ao parque de iluminação pública considerado no Quadro de Iluminação Pública – QIP. Projetado tal excedente para o período de 120 (cento e vinte) meses, equivalente a 10 (dez) anos, o montante total indevidamente faturado perfaz 50.774,40 kWh. Nos termos do art. 323 da Resolução Normativa ANEEL nº 1.000/2021 e do art. 42, parágrafo único, do Código de Defesa do Consumidor, impõe-se a repetição do indébito em dobro, totalizando 101.548,80 kWh, acrescida de atualização monetária e juros moratórios, a serem apurados até a efetiva restituição.

\item
No que concerne à interpretação do art. 473 da Resolução Normativa ANEEL nº 1.000/2021 — dispositivo que reproduz, em essência, o teor do art. 25 da revogada REN nº 414/2010 —, a ANEEL firmou entendimento no sentido de que as perdas de reatores, quando excepcionalmente admitidas, devem ser calculadas exclusivamente com base nas perdas máximas “a frio” previstas nas normas técnicas da ABNT, sendo expressamente vedadas estimativas “a quente”,percentuais genéricos ou quaisquer acréscimos destituídos de amparo tecniconormativo. Tal orientação encontra-se formalizada, entre outros atos, na Nota Técnica nº 76/2023-SMA/ANEEL e no Despacho nº 1.456/2023 STD/ANEEL, consolidando o entendimento regulatório aplicável ao caso concreto.
\end{enumerate}

\section{DOS PEDIDOS}

\begin{enumerate}[resume]

\item
Ante o exposto, requer-se à Concessionária: 

\begin{enumerate}
    \item que, em estrita observância às normas da Aneel e às cláusulas do contrato de concessão, promova a análise, solução e resposta conclusiva à presente reclamação no prazo máximo de 5 (cinco) dias úteis, nos termos do art. 408, I, da REN nº 1.000/2021;
    \item  que efetue a devolução integral dos valores indevidamente cobrados a título de “perdas em reatores” referentes a lâmpadas fluorescentes, relativamente aos últimos 10 (dez) anos, contados do primeiro protocolo da reclamação, conforme o Despacho nº 2.006/2024 (aplicação do art. 205 do Código Civil), em dobro (art. 323 da REN nº 1.000/2021 e art. 42, parágrafo único, do CDC), acrescidos de atualização monetária e juros até a efetiva devolução; 
    \item que a devolução se dê por depósito em conta corrente de titularidade do Município, na forma do art. 323, § 7º, II, da REN nº 1.000/2021, comunicando-se o cumprimento por escrito ao representante que a esta subscreve; 
    \item que, em observância aos princípios do contraditório e da ampla defesa (CF/88, art.5º, LV), bem como da eficiência e celeridade (Lei nº 9.784/1999, art. 2º), e às regras de representação do consumidor perante a distribuidora/Aneel (REN nº 1.000/2021, art. 9º), todas as comunicações, intimações e notificações relativas a esta reclamação sejam dirigidas exclusivamente ao representante legal do Município que subscreve, referencialmente pelo e-mail \emailEmpresa{} , reputando-se ineficazes aquelas enviadas a terceiros ou ao antigo representante, por risco de cerceamento de defesa;
    \item caso a Concessionária alegue a existência de dívida líquida e exigível em face do Município, que, na mesma resposta, apresente quadro analítico e documental claro e segregado, em estrita observância aos deveres de transparência e informação (CDC, art. 6º, III, e Lei nº 9.784/1999, arts. 2º e 50), contendo, obrigatoriamente: 
    \begin{itemize}
        \item débitos de consumo; 
        \item débitos oriundos de TOI de Censo;
        \item débitos oriundos de TOI comum; 
        \item débitos de faturas desagrupadas; 
        \item dívidas prescritas (dívida com mais de 60 meses). Deverá constar, ainda, a data da última atualização da dívida informada 
    \end{itemize}
\end{enumerate}


\end{enumerate}